\documentclass[11pt]{article}
\usepackage[a4paper,margin=1.85cm]{geometry}
\usepackage[T1]{fontenc}
\usepackage{lmodern}
\usepackage{amsmath,amssymb,amsthm,mathtools}
\usepackage{graphicx}
\usepackage{xcolor}
\usepackage{microtype}
\usepackage{enumitem}
\usepackage[hidelinks]{hyperref}

\definecolor{ink}{HTML}{172033}
\definecolor{accent}{HTML}{7B3F32}
\definecolor{soft}{HTML}{FAF0ED}
\color{ink}
\newtheorem{theorem}{Theorem}
\newtheorem{remark}{Remark}

\title{Transported Galerkin projections need not be uniformly bounded}
\author{Counterexample packet for Remark 7.5 of arXiv:2105.07908}
\date{11 August 2026}

\begin{document}
\maketitle

\begin{center}
\fcolorbox{accent}{soft}{\parbox{0.91\textwidth}{
\textbf{Negative answer in the abstract evolving-space framework.}
There is a smooth compatible family of a fixed Gelfand triple
$X\hookrightarrow H\hookrightarrow X^*$ and a transported orthogonal Schauder basis for which, at every $t>0$, the $H$-orthogonal projections $P_n^t$ onto the first $n$ transported vectors satisfy
\[
  \sup_n\|P_n^t|_X\|_{X\to X}=\infty.
\]
Thus the uniform $X(t)$-boundedness needed by the duality argument in Remark 7.5 is false without an additional structural assumption.}}
\end{center}

\section*{The source question and its intended reading}
In the Faedo--Galerkin construction of Alphonse, Caetano, Djurdjevac, and Elliott, an orthogonal basis $(w_j^0)$ of $X_0$ is transported to $w_j^t=\phi_t w_j^0$, and $P_n^t:H(t)\to V_n(t)$ is the $H(t)$-orthogonal projection onto
\[
 V_n(t)=\operatorname{span}\{w_1^t,\ldots,w_n^t\}.
\]
Remark 7.5 asks whether these projections are bounded uniformly in $n$ for $t>0$ \cite{Source}.

\begin{center}
\includegraphics[width=0.98\textwidth]{figures/open_problem_crop.png}
\end{center}

Literally, the displayed map in the remark has domain $V_n(t)$, on which $P_n^t$ is the identity.  The surrounding duality calculation makes the nontrivial intended question unambiguous: whether the restrictions
\[
 P_n^t|_{X(t)}:X(t)\longrightarrow V_n(t)\subset X(t)
\]
have operator norms bounded independently of $n$.  We answer that question negatively.

\section*{The compatible evolving triple}
Let $H=\ell_2(\mathbb N)$ with its standard orthonormal basis $(e_j)$.  Define weights
\[
 d_{2k-1}=k,\qquad d_{2k}=1\qquad(k\geq1),
\]
and let
\[
 X=\left\{x=(x_j):\|x\|_X^2:=\sum_{j=1}^{\infty}d_j^2|x_j|^2<\infty\right\}.
\]
Then $X$ is a separable reflexive Hilbert space, $X\hookrightarrow H$ continuously and densely, and
\[
 X\hookrightarrow H\cong H^*\hookrightarrow X^*
\]
is a Gelfand triple.  The standard basis is orthogonal in both $X$ and $H$ and is a Schauder basis of $X$; its initial coordinate projections have $X$-norm one.

Define $K$ blockwise by
\[
 Ke_{2k-1}=e_{2k},\qquad Ke_{2k}=0.
\]
The operator $K$ has norm at most one on both $H$ and $X$, and $K^2=0$.  On $0\leq t\leq1$, put
\[
 a(t)=\frac t2,\qquad \phi_t=I+a(t)K.
\]
Then
\[
 \phi_t^{-1}=I-a(t)K,
\]
so $(\phi_t)$ and $(\phi_t^{-1})$ are uniformly bounded on both $X$ and $H$, and depend smoothly on $t$ in operator norm.

Take the spaces themselves to be fixed:
\[
 X(t)=X,\qquad H(t)=H,
\]
with the standard norms, but use $\phi_t$ as the compatible evolution map.  This is allowed by the paper's definition of a compatible pair: $\phi_0=I$, each $\phi_t$ is a uniformly bounded isomorphism on both levels, and $t\mapsto\|\phi_tu\|$ is continuous.  The dual family is transported by $(\phi_t^{-1})^*$.  Moreover, all scalar products
\[
 (\phi_tu,\phi_tv)_H
\]
are polynomial in $t$ with uniformly bounded derivative.  Hence the paper's differentiability assumptions for the Gelfand-triple transport formula hold.  The evolving-space equivalence is also immediate here from the ordinary product rule for the $C^1$ families $\phi_t$ and $\phi_t^{-1}$.

Finally set
\[
 w_j^0=e_j,\qquad w_j^t=\phi_t e_j.
\]
This is exactly the paper's transported-basis construction; in particular its material derivative is zero because $\phi_t^{-1}w_j^t=e_j$ is constant.

\section*{The projection norms diverge}
Fix $t>0$ and abbreviate $a=a(t)>0$.  For every $k$,
\[
 w_{2k-1}^t=e_{2k-1}+ae_{2k},qquad w_{2k}^t=e_{2k}.
\]
The space $V_{2k-1}(t)$ contains both coordinate vectors in every earlier two-dimensional block, and in the $k$th block it contains only
\[
 f_k=e_{2k-1}+ae_{2k}.
\]
Consequently the $H$-orthogonal projection of $e_{2k}$ onto $V_{2k-1}(t)$ is simply its projection onto the line spanned by $f_k$:
\begin{equation}\label{eq:projection}
 P_{2k-1}^t e_{2k}
 =\frac{(e_{2k},f_k)_H}{\|f_k\|_H^2}f_k
 =\frac{a}{1+a^2}(e_{2k-1}+ae_{2k}).
\end{equation}
Since $\|e_{2k}\|_X=1$, while $d_{2k-1}=k$, formula \eqref{eq:projection} gives
\[
 \|P_{2k-1}^t|_X\|_{X\to X}
 \geq\|P_{2k-1}^te_{2k}\|_X
 =\frac{a}{1+a^2}\sqrt{k^2+a^2}
 \xrightarrow[k\to\infty]{}\infty.
\]

\begin{theorem}
Under the abstract hypotheses used to define the transported Galerkin spaces in Section 7.1 of the source, uniform $X(t)$-boundedness of the $H(t)$-orthogonal projections $P_n^t$ does not follow.  It can fail for every $t>0$, even when all spaces are fixed Hilbert spaces and the evolution is a smooth family of rank-one-per-block shears.
\end{theorem}

At $t=0$, one has $a(0)=0$ and $P_n^0$ is the standard coordinate projection, so $\|P_n^0\|_{X\to X}=1$.  The example therefore reproduces exactly the contrast highlighted in Remark 7.5: boundedness at the reference time does not propagate under an arbitrary compatible evolution.

\section*{Deep upgrade checks and scope}
The construction was stress-tested in three directions.
\begin{enumerate}[leftmargin=2em]
\item A unitary evolution cannot work: then $P_n^t=\phi_tP_n^0\phi_t^{-1}$ and uniform boundedness is preserved.  The shear is deliberately bounded and invertible on $X$ and $H$ but has an adjoint that interacts badly with the alternating weights.
\item A single rank-one shear produces only one potentially bad partial projection.  The block-diagonal nilpotent shear above repeats the obstruction at every odd index while retaining uniform operator bounds and the explicit inverse $I-aK$.
\item The example is abstract, not geometric.  It does not rule out uniform boundedness for Laplace eigenbases transported by smooth domain diffeomorphisms or for other specially structured PDE bases.  It proves that such a conclusion requires an assumption beyond compatibility and smooth transport.
\end{enumerate}

\section*{Bounded novelty check}
On 11 August 2026, the exact arXiv id, title, Remark 7.5 phrasing, transported Galerkin projections, evolving Banach spaces, and uniform projection boundedness were checked in the run indexes and by bounded web/arXiv search.  The searches found the source and general literature on Galerkin and uniformly bounded projections, but no answer to this exact remark.  This is evidence for novelty, not a definitive guarantee.

\begin{thebibliography}{9}
\bibitem{Source}
A. Alphonse, D. Caetano, A. Djurdjevac, and C. M. Elliott,
\emph{Function spaces, time derivatives and compactness for evolving families of Banach spaces with applications to PDEs},
J. Differential Equations \textbf{353} (2023), 268--338; arXiv:2105.07908.
\end{thebibliography}

\end{document}
